\section{Assume that, the population of Rajshahi city increases at a rate proportional to the number of inhabitants at any time $t$. If the population in double in 40 years, how many years it will be triple.}
According to the problem, we get:
\begin{align*}
    \frac{dN}{dt} &= kN \sidenote{1}
\end{align*}

The general solution of (1) is:
\begin{align*}
    N(t) &= N_0 e^{kt} \sidenote{2}
\end{align*}

Given that when $t = 40$, $N = 2N_0$:
\begin{align*}
    2N_0 &= N_0 e^{40k} \\
    40k &= \ln 2 \\
    k &= \frac{\ln 2}{40} \\
    \implies k &\approx 0.017328
\end{align*}

Suppose the population will triple in time $T$, i.e., $N(T) = 3N_0$:
\begin{align*}
    3N_0 &= N_0 e^{kT} \\
    e^{kT} &= 3 \\
    kT &= \ln 3 \\
    T &= \frac{\ln 3}{k} \\
    T &= \frac{\ln 3}{0.017328} \\
    \implies T &\approx 63.400986 \text{ years}
\end{align*}

So, the population will triple in $63.400986$ years (or approximately $63.40$ years from $t = 0$).